\chapter*{Resumen}
\addcontentsline{toc}{chapter}{Resumen}
\markboth{Resumen}{Resumen}

La enseñanza del Método de los Elementos Finitos (MEF) enfrenta una brecha pedagógica persistente: los textos clásicos exponen la formulación matricial con rigor pero en abstracto, y el software profesional opera como una caja negra que oculta el procedimiento que el estudiante necesita comprender. Esta tesis aborda dicha brecha mediante el desarrollo de EduFEM, una herramienta educativa interactiva en español para el análisis estructural de problemas bidimensionales de elasticidad lineal en tensión plana y deformación plana, entendido como el análisis tenso-deformacional de medios continuos. La herramienta combina un motor de cálculo con elementos isoparamétricos cuadriláteros Q4 y Q9, una interfaz gráfica que integra pre-proceso, proceso y post-proceso sobre un único lienzo de malla, y ocho módulos educativos que exponen paso a paso cada etapa del cálculo, del mapeo isoparamétrico al ensamblaje global, sobre el modelo del alumno. La corrección del motor se sometió a una batería de verificación y validación. Con el método de soluciones manufacturadas, en ambos estados planos y sobre mallas regulares y distorsionadas, se confirmaron las tasas de convergencia asintóticas teóricas del desplazamiento —orden dos en norma $L^2$ y orden uno en seminorma $H^1$ para Q4; tres y dos para Q9— y la convergencia del campo de tensiones recuperado. La validación —contraste con una solución analítica y con un modelo comercial, no con ensayos físicos— arrojó en la viga de Timoshenko errores del 0,04\,\% en la tensión normal, del 0,26\,\% en la flecha y de hasta el 2,9\,\% en la tensión cortante, y diferencias del 0,21\,\% y de hasta el 0,56\,\% frente a SAP2000. La membrana de Cook evidenció el bloqueo por cortante (\emph{shear-locking}) del Q4 frente a la convergencia rápida del Q9, que a igual número de grados de libertad (GDL) reduce el error en un orden de magnitud. El aporte principal es EduFEM, una plataforma que articula teoría, implementación y experimentación numérica para el aprendizaje del MEF, con una memoria de cálculo automática, paso a paso, que el estudiante puede replicar y verificar a mano.

\noindent\textbf{Palabras clave:} método de elementos finitos, software educativo, elementos isoparamétricos Q4/Q9, elasticidad plana, verificación y validación, memoria de cálculo